% Eigenständiger Prosabeweis ohne lokale nummerierte Deklarationen.
\section*{Fragestellung und Beispiel}

Wir beginnen mit zwei verschiedenen Zuständen \(a\) und \(b\).
Der Anfangszustand sei \(a\); die Übergangsregel schickt \(a\)
nach \(b\) und lässt \(b\) fest. In Zeichen:
\[
  M=\{a,b\},\qquad m_0=a,\qquad f(a)=b,\quad f(b)=b.
\]
Die Regel \(f\colon M\to M\) ist damit vollständig gegeben.
Gesucht ist eine Funktion \(F\colon\mathbb N\to M\), die jeder
natürlichen Stelle den dort erreichten Zustand zuordnet. Dabei
bezeichnet \(\Succ(n)\) den Nachfolger von \(n\). Die ersten
Werte sind durch Anfangszustand und Übergangsregel festgelegt:
\[
  F(0)=a,\qquad F(\Succ(0))=b,\qquad
  F(\Succ(\Succ(0)))=b.
\]
Auch jeder weitere Übergang von \(b\) führt wieder nach \(b\).
Die natürlichen Stellen bleiben verschieden, obwohl sich ihre
zugeordneten Werte wiederholen dürfen.

\DedekindRecursionDiagram

Was bedeutet es aber, diese Vorschrift an \emph{allen} natürlichen
Stellen auszuführen? Das Hinschreiben weiterer Werte zeigt zunächst
nur, wie eine passende Funktion aussehen müsste. Zu begründen ist,
dass eine einzige Funktion auf ganz \(\mathbb N\) existiert und
dass keine andere Funktion dieselben Vorgaben erfüllt. Genau dies
leistet der Dedekindsche Rekursionssatz für jede Menge von Zuständen
und jede vorgegebene Übergangsabbildung.

\par\Needspace{16\baselineskip}
\section*{Aussage und Beweisidee}

Eine Rekursionsvorschrift gibt einen Anfangswert und eine Regel für den
Übergang zum nächsten Wert vor. Der Dedekindsche Rekursionssatz zeigt,
dass diese Angaben eine Funktion auf allen natürlichen Zahlen eindeutig
bestimmen. Genauer: Ist \(M\) eine Menge, \(m_0\in M\) und
\(f\colon M\to M\), so gibt es genau eine Funktion
\(F\colon\mathbb N\to M\) mit
\[
  F(0)=m_0,
  \qquad
  F(\Succ(n))=f(F(n))
  \quad(n\in\mathbb N).
\]
Dies ist die Aussage von
\FormulaRefAuto{DedekindRecursionTheorem}[thm] in Band~10.
Die Abbildung \(f\) ist bereits gegeben; zu konstruieren ist die
Funktion \(F\), welche die aufeinanderfolgenden Werte erfasst.
Weder Injektivität noch Surjektivität von \(f\) werden vorausgesetzt.

Für den Beweis werden die Peano-Eigenschaften verwendet:
\(0\in\mathbb N\), jeder natürliche Nachfolger liegt wieder in
\(\mathbb N\), kein Nachfolger ist \(0\), und die
Nachfolgerabbildung ist injektiv. Hinzu kommt das Induktionsprinzip:
Gilt eine Aussage für \(0\), und folgt aus ihrer Gültigkeit für
\(n\) stets ihre Gültigkeit für \(\Succ(n)\), so gilt sie für
alle \(n\in\mathbb N\).
Addition und eine Ordnung auf \(\mathbb N\) werden nicht benötigt.

Die Gleichungen legen nahe, zunächst \(m_0\), dann \(f(m_0)\)
und danach \(f(f(m_0))\) zu betrachten. Damit ist aber die Existenz
einer einzigen Funktion auf ganz \(\mathbb N\) noch nicht bewiesen.
Wir konstruieren deshalb ihren Graphen als Menge geordneter Paare.
Dazu betrachten wir alle Teilmengen von \(\mathbb N\times M\),
die das Startpaar enthalten und unter dem gewünschten Übergang
abgeschlossen sind. Ihr Schnitt ist die kleinste solche Menge.

Diese Schnittmenge heißt \emph{Rekursionskern}. Dass sie bei jedem
\(n\) genau einen Wert enthält, ist der entscheidende Schritt.
Eine unter der Übergangsregel abgeschlossene Relation kann zunächst
mehrere Werte an derselben Stelle enthalten. Auch beim kleinsten
solchen Objekt dürfen wir die Funktionseigenschaft daher noch nicht
voraussetzen.
Der Rekursionsabschluss liefert jeweils einen Wert in der
Nachfolgerstufe. Geeignete Fixierungsmengen schließen weitere Werte
aus. Erst danach lesen wir den Kern als Funktion und weisen die
Eindeutigkeit unter allen Funktionen mit den vorgegebenen
Rekursionsgleichungen nach.

Die Fixierungsmengen nutzen die Minimalität auf folgende Weise:
Wir beschränken
die an einer bestimmten Stelle zugelassenen Werte und zeigen, dass
die eingeschränkte Menge die Abschlussbedingungen weiterhin erfüllt.
Dann muss sie den kleinsten Kern enthalten. Die ausgeschlossenen
Paare können folglich schon vorher nicht zum Kern gehört haben.
So gewinnt man aus der Minimalität die Eindeutigkeit der Werte.

\par\Needspace{9\baselineskip}
\section*{Konstruktion und Beweis}
\begin{proof}[Beweis]
\textbf{Der Rekursionskern.}

Wir halten \(M\), \(m_0\in M\) und \(f\colon M\to M\)
für den gesamten Beweis fest. Eine Menge \(H\) heißt bezüglich
dieser Daten \emph{rekursionsadmissibel}, wenn sie die folgenden
drei Bedingungen erfüllt:
\[
  \begin{aligned}
    &H\subseteq\mathbb N\times M,\\
    &(0,m_0)\in H,\\
    &\forall n\in\mathbb N\,\forall a\in M\,
      \bigl((n,a)\in H
        \ \Longrightarrow\ (\Succ(n),f(a))\in H\bigr).
  \end{aligned}
\]
Wir schreiben dafür \(\AdmRec{m_0}{f}(H)\).
Admissibilität fordert zunächst keine Eindeutigkeit des Werts
\(a\) zu einem gegebenen \(n\). Es dürfen also verschiedene
Paare mit derselben ersten Koordinate in \(H\) liegen.
Im Eingangsbeispiel erfüllt etwa die gesamte Relation
\(\mathbb N\times\{a,b\}\) die Bedingungen. Sie enthält sowohl
\((0,a)\) als auch \((0,b)\), obwohl der Anfangswert nur \(a\)
sein soll. Die Abschlussbedingungen allein reichen also noch nicht
aus, um einen Funktionsgraphen auszuzeichnen.

Der volle Produktraum \(\mathbb N\times M\) ist admissibel:
Das Startpaar liegt darin, weil \(0\in\mathbb N\) und
\(m_0\in M\). Gehört \((n,a)\) zum Produktraum, so liegen
\(\Succ(n)\) in \(\mathbb N\) und \(f(a)\) in \(M\).
Also liegt auch \((\Succ(n),f(a))\) darin.
Somit ist die Menge
\[
  \mathcal A
    :=\{H\in\powerset(\mathbb N\times M)
         \mid \AdmRec{m_0}{f}(H)\}
\]
nichtleer. Als Teilmenge einer Potenzmenge ist \(\mathcal A\)
eine Menge. Wir können ihren Schnitt bilden und setzen
\[
  K:=\RecCore{m_0}{f}:=\bigcap_{H\in\mathcal A}H.
\]
Ein Paar liegt genau dann in \(K\), wenn es in jeder admissiblen
Menge liegt.

Der Kern ist selbst admissibel. Zunächst gilt
\(K\subseteq\mathbb N\times M\), weil der Produktraum ein
Mitglied von \(\mathcal A\) ist. Ferner liegt \((0,m_0)\)
in jedem \(H\in\mathcal A\), also auch in deren Schnitt.
Ist schließlich \((n,a)\in K\), so liegt dieses Paar in jedem
\(H\in\mathcal A\). Dessen Rekursionsabschluss liefert
\((\Succ(n),f(a))\in H\). Da dies für jedes solche \(H\)
gilt, folgt
\[
  (n,a)\in K
  \quad\Longrightarrow\quad
  (\Succ(n),f(a))\in K.
\]

Außerdem ist \(K\) in jeder admissiblen Menge enthalten:
\[
  \AdmRec{m_0}{f}(H)\quad\Longrightarrow\quad K\subseteq H.
\]
Zusammen mit seiner eigenen Admissibilität macht dies \(K\) zur
\emph{kleinsten} admissiblen Menge bezüglich der Inklusion.
Diese Minimalität wird nun benutzt, um die Werte des Kerns
festzulegen. In den Beweistabellen ist sie als
\FormulaRefAuto{RecCoreLeastAdmissible}[thm] festgehalten.

\par\Needspace{7\baselineskip}
\medskip\noindent\textbf{Genau ein Wert in jeder Stufe.}

Für ein \(n\in\mathbb N\) nennen wir die Paare \((n,a)\in K\)
die \emph{Stufe bei \(n\)}. Wir zeigen durch Induktion
\[
  U(n)\;:\Longleftrightarrow\;
  \exists!a\in M\,\bigl((n,a)\in K\bigr).
\]
Existenz und Eindeutigkeit werden dabei gemeinsam bewiesen.

\medskip\noindent
\textbf{Die Nullstufe.}
Die Existenz ist bereits bekannt: \((0,m_0)\in K\).
Um andere Werte bei \(0\) auszuschließen, betrachten wir
\[
  Z:=\{(n,a)\in\mathbb N\times M
         \mid n=0\ \Longrightarrow\ a=m_0\}.
\]
Die Menge \(Z\) enthält im Produktraum alle Paare, deren erste
Koordinate von \(0\) verschieden ist; bei \(0\) lässt sie nur
das Paar \((0,m_0)\) zu.

Die Menge \(Z\) ist admissibel. Sie ist eine Teilmenge des
Produktraums und enthält das Startpaar. Ist \((n,a)\in Z\),
so liegt \((\Succ(n),f(a))\) im Produktraum. Weil
\(\Succ(n)\neq0\), erfüllt dieses Nachfolgerpaar auch die
zusätzliche Bedingung in der Definition von \(Z\).
Es gehört somit zu \(Z\).

Aus der Minimalität des Kerns folgt \(K\subseteq Z\).
Für jedes \((0,a)\in K\) gilt daher \((0,a)\in Z\) und
folglich \(a=m_0\). Zusammen mit dem vorhandenen Startpaar
ergibt das \(U(0)\).
Im Eingangsbeispiel schließt dieser Schritt insbesondere das
zusätzliche Paar \((0,b)\) aus, das im vollen Produktraum
noch enthalten war.

\medskip\noindent
\textbf{Die Nachfolgerstufe.}
Sei \(x\in\mathbb N\), und es gelte \(U(x)\).
Dann gibt es ein \(u\in M\) mit
\[
  (x,u)\in K,
  \qquad
  \forall c\in M\,
    \bigl((x,c)\in K\ \Longrightarrow\ c=u\bigr).
\]
Der Rekursionsabschluss von \(K\) liefert das Paar
\((\Succ(x),f(u))\in K\). Sein zweiter Eintrag liegt in
\(M\), weil \(f\colon M\to M\). Damit besitzt die Stufe
bei \(\Succ(x)\) mindestens einen Wert.

Zur Eindeutigkeit definieren wir die Fixierungsmenge
\[
  \Phi_{x,u}
    :=\{(y,d)\in K
          \mid y=\Succ(x)\ \Longrightarrow\ d=f(u)\}.
\]
Diese Bedingung würde nur solche Paare der Stufe \(\Succ(x)\)
ausschließen, deren zweiter Eintrag nicht \(f(u)\) ist.
An allen anderen Stufen bleiben die Paare des Kerns erhalten.
Wir zeigen, dass diese eingeschränkte Menge weiterhin admissibel
ist. Dann zwingt die Minimalität \(K\) dazu, schon in
\(\Phi_{x,u}\) enthalten zu sein.

Zunächst gilt
\(\Phi_{x,u}\subseteq K\subseteq\mathbb N\times M\).
Auch das Startpaar bleibt erhalten: Es liegt in \(K\), und
\(0\neq\Succ(x)\). Die Bedingung
\(0=\Succ(x)\Longrightarrow m_0=f(u)\) ist daher erfüllt.
Also ist \((0,m_0)\in\Phi_{x,u}\).

Für den Rekursionsabschluss sei \((y,d)\in\Phi_{x,u}\).
Dann ist \((y,d)\in K\); insbesondere liegen \(y\) in
\(\mathbb N\) und \(d\) in \(M\). Der Abschluss von
\(K\) gibt
\[
  (\Succ(y),f(d))\in K.
\]
Für die Zugehörigkeit zu \(\Phi_{x,u}\) müssen wir zusätzlich
die Implikation
\[
  \Succ(y)=\Succ(x)\quad\Longrightarrow\quad f(d)=f(u)
\]
nachweisen. Nehmen wir dazu \(\Succ(y)=\Succ(x)\) an.
Die Injektivität der Nachfolgerabbildung liefert \(y=x\).
Aus \((y,d)\in K\) folgt somit \((x,d)\in K\).
Wegen der Eindeutigkeit in der Stufe bei \(x\) ist \(d=u\),
also \(f(d)=f(u)\). Die erforderliche Implikation ist bewiesen,
und damit liegt \((\Succ(y),f(d))\) in \(\Phi_{x,u}\).

Folglich ist \(\Phi_{x,u}\) admissibel. Aus der Minimalität
des Kerns folgt \(K\subseteq\Phi_{x,u}\). Nach ihrer Definition
gilt zugleich \(\Phi_{x,u}\subseteq K\), also
\(\Phi_{x,u}=K\). Die Einschränkung hat somit tatsächlich kein
Paar des Kerns entfernt.
Ist nun \((\Succ(x),b)\in K\), so liegt dieses Paar auch
in \(\Phi_{x,u}\), und deren definierende Bedingung ergibt
\(b=f(u)\). Jeder Wert der Nachfolgerstufe ist also der bereits
gefundene Wert \(f(u)\). Damit gilt \(U(\Succ(x))\).

\DedekindFixingDiagram

Das Induktionsprinzip liefert insgesamt
\[
  \forall n\in\mathbb N\,
    \exists!a\in M\,\bigl((n,a)\in K\bigr).
\]
An jeder natürlichen Stelle ist der Kern somit vorhanden und
eindeutig. Im Nachfolgerschritt wurde nur der eindeutige Wert an
einer fest gewählten Stelle verwendet; eine gleichzeitige Auswahl
von Werten an allen Stellen war nicht nötig.
Die eindeutige Existenz an jeder Stelle ist die Aussage von
\FormulaRefAuto{RecCoreStageUnique}[thm] in den Beweistabellen.

\par\Needspace{6\baselineskip}
\medskip\noindent\textbf{Die Funktion und ihre Eindeutigkeit.}

Eine Teilmenge von \(\mathbb N\times M\), die jedem
\(n\in\mathbb N\) genau ein \(a\in M\) zuordnet, ist
der Graph einer Funktion \(\mathbb N\to M\).
Diese Voraussetzungen haben wir für \(K\) nachgewiesen.
Wir können daher den Kern selbst als Funktion lesen und setzen
\(F:=K\). Dabei bedeutet \(F(n)=a\) genau, dass
\((n,a)\in K\).

Das Startpaar ergibt sofort \(F(0)=m_0\).
Für ein beliebiges \(n\in\mathbb N\) gilt
\((n,F(n))\in K\). Der Rekursionsabschluss liefert
\[
  (\Succ(n),f(F(n)))\in K.
\]
Weil \(K\) der Graph von \(F\) ist, folgt
\[
  F(\Succ(n))=f(F(n)).
\]
Damit ist die Existenz einer Funktion mit den verlangten
Eigenschaften bewiesen.

Zur Eindeutigkeit sei \(G\colon\mathbb N\to M\) eine weitere
Funktion mit
\[
  G(0)=m_0,
  \qquad
  G(\Succ(n))=f(G(n))
  \quad(n\in\mathbb N).
\]
Wir beweisen \(G(n)=F(n)\) durch Induktion.
Bei \(0\) stimmen beide Funktionen mit \(m_0\) überein.
Gilt \(G(n)=F(n)\), so ergibt sich
\[
  \begin{aligned}
    G(\Succ(n))
      &=f(G(n))\\
      &=f(F(n))\\
      &=F(\Succ(n)).
  \end{aligned}
\]
Also stimmen \(G\) und \(F\) an jeder natürlichen Stelle
überein. Da sie denselben Definitionsbereich haben, sind sie
als Funktionen gleich. Damit ist der Dedekindsche Rekursionssatz
vollständig bewiesen.
\end{proof}

\par\Needspace{38\baselineskip}
\section*{Was die Konstruktion sichtbar macht}

Im Beweis treten zwei verschiedene Eindeutigkeitsaussagen auf.
Zuerst wird innerhalb der Relation \(K\) für jede natürliche
Stelle genau ein Wert nachgewiesen. Erst dadurch ist \(K\)
eine Funktion. Anschließend wird diese Funktion mit jeder anderen
Funktion verglichen, die denselben Anfangswert und dieselbe
Übergangsregel besitzt. Die zweite Induktion zeigt, dass auch
die gesamte Lösung eindeutig bestimmt ist.

Induktion und Rekursion erfüllen dabei unterschiedliche Aufgaben.
Mit dem Induktionsprinzip weisen wir Aussagen an allen natürlichen
Stellen nach. Der Rekursionssatz sichert die Existenz und
Eindeutigkeit einer Funktion mit den vorgeschriebenen Werten.
In unserem Beweis greifen beide zusammen: Die Schnittkonstruktion
liefert eine Relation, und Induktion weist deren Funktionseigenschaft
sowie die Eindeutigkeit der Lösung nach.

Das Eingangsbeispiel zeigt außerdem, dass Eindeutigkeit eines Werts
keine Injektivität bedeutet. An verschiedenen Stellen darf derselbe
Wert stehen: Hier erhalten alle Nachfolger den Wert \(b\),
während an der Stelle \(0\) der Wert \(a\) steht. Weder die
gegebene Übergangsabbildung \(f\) noch die daraus gewonnene
Funktion \(F\) muss injektiv sein.

Im Hauptwerk bezeichnet \(\RecFun{m_0}{f}\) die durch den Satz
eindeutig bestimmte Funktion
(\FormulaRefAuto{RecFunDef}[def]). Ihre Definition benötigt nur
den Satz und seine Rekursionsgleichungen. Der hier konstruierte Kern
erfüllt diese Gleichungen und ist deshalb genau diese Funktion:
\[
  \RecFun{m_0}{f}=\RecCore{m_0}{f}.
\]
Der Satz rechtfertigt so die späteren rekursiven Definitionen in
Band~10, insbesondere den anschließenden Aufbau der Addition und
Multiplikation. Diese Operationen werden im vorstehenden Beweis
noch nicht vorausgesetzt. Wie beim Lesebeweis zu Cantor--Bernstein
wird ein kleinstes abgeschlossenes Objekt durch einen Schnitt
gewonnen; hier ist dieses Objekt ein Graph von Zahlen und Werten.
Die Definitionen der Hilfsmengen und sämtliche formalen
Ableitungen stehen in den zugehörigen \DedekindProofsLink{}.
